
\section{GPU chaining}
\label{sec:chaining-b2}

\subsection{G3SA algorithm}

The G3SA paper divides GPU kernel into 4 stages:
\begin{enumerate}
\item \textbf{First forward-scoring pass:} Either use short kernel (specialized for short sequences) or tiled kernel (tiled kernel is just dividing 
sequence into smaller chunks and applying modified short kernel). This stage computes total scores as normal chaining $f_j + E(j,i)$.
\item \textbf{First backtracking and packing:}  G3SA applies minimap2's
    cut and marking rule, reverses accepted paths, and packs their anchors into a buffer.
\item \textbf{Conditional ranking:} If the difference between sizes of query and target is big (see Figure~\ref{fig:chaining-pipelines}), then 
    kernel compares simplified\(h(j)\) and calculates one complete link score for the best canidate (similar to how minimap2 does RMQ rescoring mentioned earlier)
\item \textbf{Final backtracking:} This is done via minimap2 scalar backtracking algorithm
\end{enumerate}

\info{G3SA AVL code}{The repository contains CUDA AVL support code and a
commented experimental entry point.  The active chaining pipeline does not
call that entry point.  Its second pass uses array-based short and tiled
kernels.}

\begin{figure}[H]
\centering
\begin{tikzpicture}[x=1cm,y=1cm,font=\footnotesize,>=Latex]
\draw[chain muted] (.10,3.25) rectangle (1.55,4.15);
\node at (.825,3.70) {anchors};

\draw[chain blue] (2.10,2.92) rectangle (4.65,4.48);
\node[font=\footnotesize\bfseries] at (3.375,4.15)
    {first-pass score};
\draw[fill=chainBlue!6,draw=chainBlue!70!black]
    (2.30,3.08) rectangle (3.28,3.72);
\node at (2.79,3.40) {short};
\draw[fill=chainBlue!6,draw=chainBlue!70!black]
    (3.47,3.08) rectangle (4.45,3.72);
\node at (3.96,3.40) {tiled};

\draw[chain green] (5.20,3.15) rectangle (7.55,4.25);
\node[align=center] at (6.375,3.70)
    {backtrack\\and pack};

\draw[chain orange] (8.10,2.45) rectangle (10.72,4.95);
\node[align=center,text width=2.28cm,inner sep=0pt] at (9.41,3.70)
    {Is the found chain vs. target size greater than 1000 or 10\% of target size?};

\draw[chain green] (11.35,3.15) rectangle (14.15,4.25);
\node[align=center] at (12.75,3.70)
    {no: keep first-pass\\chain records};

\draw[chain blue] (8.10,.75) rectangle (10.72,1.95);
\node[align=center] at (9.41,1.35)
    {pool and sort\\rank again};

\draw[chain green] (11.35,.75) rectangle (14.15,1.95);
\node[align=center] at (12.75,1.35)
    {same backtracking\\final records};

\draw[chain flow] (1.55,3.70)--(2.10,3.70);
\draw[chain flow] (4.65,3.70)--(5.20,3.70);
\draw[chain flow] (7.55,3.70)--(8.10,3.70);
\draw[chain flow] (10.72,3.70)--(11.35,3.70)
    node[midway,above] {no};
\draw[chain flow] (9.41,2.45)--(9.41,1.95)
    node[midway,right] {yes};
\draw[chain flow] (10.72,1.35)--(11.35,1.35);
\end{tikzpicture}
\caption{Simplified depictions of G3SA pipeline}
\label{fig:chaining-pipelines}
\end{figure}

\leavevmode\coderef{g3sa/minimap2/src/kernel_wrapper.cu}
{G3SA chaining pipeline}{chaining_kernel_wrapper}{775}
\leavevmode\coderef{g3sa/minimap2/src/krmq_cuda.cu}
{Commented CUDA AVL entry point}{gpu_chain_rmq}{406}

The main difference between G3SA is that it uses specialised paths depending on data size and it calculates rows instead of columns in the score matrix for anchors:
the first kernel instead of looking at previous anchors and updating itself, looks at next anchors and updates them~\ref{fig:chaining-work-split}.

\begin{figure}[H]
\centering
\begin{tikzpicture}[x=1cm,y=1cm,font=\footnotesize,>=Latex]
    \node[anchor=west,font=\bfseries] at (.15,6.75)
    {(a) minimap2 B1: fix target \(i\)};
    \node[anchor=west,font=\bfseries] at (7.35,6.75)
    {(b) G3SA short: fix source \(j\)};

    \foreach \r/\c in {
        0/1,0/2,0/3,0/4,0/5,
        1/2,1/3,1/4,1/5,
        2/3,2/4,2/5,
        3/4,3/5,
        4/5}
    \draw[fill=chainGray!12,draw=chainGray!80!black]
        (1.35+.62*\c-.27,5.72-.62*\r-.27)
        rectangle
        (1.35+.62*\c+.27,5.72-.62*\r+.27);
    \foreach \r in {0,...,4}
    \draw[fill=chainBlue!20,draw=chainBlue!90!black,very thick]
        (1.35+.62*5-.27,5.72-.62*\r-.27)
        rectangle
        (1.35+.62*5+.27,5.72-.62*\r+.27);
    \draw[fill=chainOrange!24,draw=chainOrange!95!black,very thick]
    (1.35+.62*5-.27,5.72-.62*5-.27)
    rectangle
    (1.35+.62*5+.27,5.72-.62*5+.27);
    \node at (1.35+.62*5,5.72-.62*5) {\(i\)};
    \foreach \c in {0,...,5}
    \node[anchor=south] at (1.35+.62*\c,6.06) {\scriptsize \(i_{\c}\)};
    \foreach \r in {0,...,5}
    \node[anchor=east] at (1.00,5.72-.62*\r) {\scriptsize \(j_{\r}\)};
    \draw[chain flow,draw=chainBlue!90!black]
    (5.28,2.78)--(5.28,5.40);
    \node[rotate=90,anchor=south,text=chainBlue!90!black]
    at (5.43,4.08) {newest to older};

    \foreach \r/\c in {
        0/1,0/2,0/3,0/4,0/5,
        1/2,1/3,1/4,1/5,
        2/3,2/4,2/5,
        3/4,3/5,
        4/5}
    \draw[fill=chainGray!12,draw=chainGray!80!black]
        (8.55+.62*\c-.27,5.72-.62*\r-.27)
        rectangle
        (8.55+.62*\c+.27,5.72-.62*\r+.27);
    \foreach \c in {3,...,5}
    \draw[fill=chainBlue!20,draw=chainBlue!90!black,very thick]
        (8.55+.62*\c-.27,5.72-.62*2-.27)
        rectangle
        (8.55+.62*\c+.27,5.72-.62*2+.27);
    \draw[fill=chainOrange!24,draw=chainOrange!95!black,very thick]
    (8.55+.62*2-.27,5.72-.62*2-.27)
    rectangle
    (8.55+.62*2+.27,5.72-.62*2+.27);
    \node at (8.55+.62*2,5.72-.62*2) {\(j\)};
    \foreach \c in {0,...,5}
    \node[anchor=south] at (8.55+.62*\c,6.06) {\scriptsize \(i_{\c}\)};
    \foreach \r in {0,...,5}
    \node[anchor=east] at (8.20,5.72-.62*\r) {\scriptsize \(j_{\r}\)};
    \draw[chain flow,draw=chainBlue!90!black]
    (12.00,4.48)--(13.20,4.48);
    \node[anchor=south,text=chainBlue!90!black] at (12.60,4.58)
    {update \(i>j\)};

    \draw[fill=chainBlue!12,draw=chainBlue!90!black]
    (.80,.45) rectangle (3.75,1.25);
    \node at (2.275,.85) {read earlier \(j\)};
    \draw[fill=chainOrange!18,draw=chainOrange!95!black]
    (4.55,.45) rectangle (6.30,1.25);
    \node at (5.425,.85) {one \(i\)};
    \draw[chain flow] (3.75,.85)--(4.55,.85);

    \draw[fill=chainOrange!18,draw=chainOrange!95!black]
    (7.55,.45) rectangle (9.30,1.25);
    \node at (8.425,.85) {one \(j\)};
    \draw[fill=chainBlue!12,draw=chainBlue!90!black]
    (10.10,.45) rectangle (13.55,1.25);
    \node at (11.825,.85) {write later \(i\)};
    \draw[chain flow] (9.30,.85)--(10.10,.85);
\end{tikzpicture}
\caption{Different approaches (transposed) of minimap2 and G3SA}
\label{fig:chaining-work-split}
\end{figure}

After long analysis I decided to implement the simpler approach. Instead of writing multiple kernels, just naively port existing chaining logic to Vulkan API.
\leavevmode\coderef{g3sa/minimap2/src/chaining_kernel.cu}
  {G3SA short ranking kernel}{mm_compute_chain_short_rmq}{970}
\leavevmode\coderef{g3sa/minimap2/src/kernel_wrapper.cu}
  {G3SA ranking-pass launches}{chaining_kernel_wrapper_rmq}{595}

\subsection{BioseqDB Vulkan tiled DP}
\label{sec:chaining-vulkan}

BioseqDB's \texttt{vulkan\_tiled\_dynamic\_programming} is the usual chaining algorithm that fixes one anchor \(i\) and examines earlier predecessors in
batches of 32.

The host first sets
\[
  D=\max\{\texttt{max\_gap},\texttt{bandwidth}\}.
\]
and sorts anchors to filter them by adjacent reference gap greater than \(D\).
For a fixed target, lane \(u\) first scores predecessor \(i-u-1\).  The next
iteration moves all lanes 32 positions to the back.  Every active lane
calculates the complete value \(f_j+E(j,i)\) to then select one
pair \((\text{score},j)\).

\begin{algorithm}[H]
\caption{BioseqDB Vulkan forward chaining for one hard range}
\label{alg:chaining-vulkan}
\footnotesize
\begin{algorithmic}[1]
  \Require sorted range \(a_b,\ldots,a_{e-1}\), distance \(D\),
    scan limit \(K\), skip parameter \(S\)
  \Ensure scores \(f\) and absolute predecessor indices \(p\)
  \State \(U\gets\infty\) if \(S\leq0\), otherwise
    \(U\gets\lfloor S/32\rfloor+1\)
  \For{\(i\gets b\) \textbf{to} \(e-1\)}
    \State \(f_i\gets\ell_i,\quad p_i\gets-1\)
    \State \(o\gets1,\quad u\gets0\)
    \While{\(o\leq K\land i-o\geq b\)}
      \State \(B\gets\) the longest prefix of
        \((i-o,\ldots,i-o-31)\) with \(j\geq b\),
        \(i-j\leq K\), and \(r_i-r_j\leq D\)
      \If{\(B=\varnothing\)}
        \State \textbf{break}
      \EndIf
      \ForAll{\(j\in B\) in 32 lanes}
        \State \(c_j\gets f_j+E(j,i)\)
            \Comment{invalid links use \(-\infty\)}
      \EndFor
      \State \(j^\star\gets
        \operatorname*{arg\,max}_{j\in B}(c_j,j)\)
            \Comment{larger \(j\) breaks a score tie}
      \State \(c\gets c_{j^\star}\)
      \If{\(c>f_i\)}
        \State \(f_i\gets c,\quad p_i\gets j^\star,\quad u\gets0\)
      \ElsIf{\(S>0\)}
        \State \(u\gets u+1\)
        \If{\(u\geq U\)}
          \State \textbf{break}
        \EndIf
      \EndIf
      \If{\(|B|<32\)}
        \State \textbf{break}
      \EndIf
      \State \(o\gets o+32\)
    \EndWhile
    \State \(\operatorname{await}(\mathrm{targetUpdates}_i)\)
          \Comment{make \(f_i,p_i\) visible before target \(i+1\)}
  \EndFor
\end{algorithmic}
\end{algorithm}
